%%=============================================================================
%% Methodologie
%%=============================================================================

\chapter{\IfLanguageName{dutch}{Methodologie}{Methodology}}%
\label{ch:methodologie}

\section{Onderzoeksopzet}

Het doel van deze bachelorproef is het ontwikkelen van een interpreteerbare en reproduceerbare fysieke spelersscore ter ondersteuning van objectieve voetbalscouting. De realisatie hiervan hanteert een data-gedreven methodologie waarbij machine learning-modellen worden ingezet om het relatieve belang van fysieke prestatie-indicatoren te bepalen.

Er wordt gekozen voor een hybride aanpak waarbij supervised learning gecombineerd wordt met model-interpretatietechnieken. Daarnaast bekijkt het onderzoek naar het belang van de fysieke prestatie-indicatoren per positiegroep, zo worden de fysieke rolvereisten per positie berekend.

\section{Dataset}

De spelersdataset bestaat uit wedstrijdgerelateerde fysieke prestatie-indicatoren per speler per wedstrijd, opgehaald uit de Skillcorner data. Alle variabelen worden genormaliseerd per 90 minuten om onderlinge vergelijkbaarheid te garanderen.

De andere dataset bestaat uit wedstrijdresultaten, waar per wedstrijd het aantal doelpunten voor de thuis- en uitploeg te vinden is.

Het onderzoek neemt ook contextvariabelen zoals speelminuten mee voor normalisatiecontrole.

\section{Positie-specifieke modellering}

Aangezien fysieke vereisten sterk verschillen per speelpositie, wordt de dataset opgesplitst per positie (bijvoorbeeld centrale verdedigers, flankverdedigers, centrale middenvelders en flankaanvallers). Voor elke positie wordt een afzonderlijk machine learning-model getraind.

Deze aanpak zorgt ervoor dat spelers uitsluitend worden vergeleken met positiegenoten, wat structurele bias in de scoreberekening voorkomt.

\section{Supervised learning}

Voor elke positie wordt een classificatiemodel ontwikkeld met als doelvariabele een proxy voor fysieke waarde binnen de ploeg. Als proxy wordt gekozen voor aantal speelminuten in volgende wedstrijd.

Er wordt gebruikgemaakt van een gradient boosting model (bijvoorbeeld XGBoost) om niet-lineaire relaties tussen fysieke variabelen en de doelvariabele te kunnen modelleren.

De dataset wordt opgesplitst in trainings- en testdata, waarbij cross-validatie wordt toegepast om overfitting te vermijden. Modelperformantie wordt geëvalueerd aan de hand van geschikte metrics, zoals accuracy, F1-score of RMSE, afhankelijk van het type doelvariabele.

\section{Modelinterpretatie via SHAP}

Om het relatieve belang van individuele fysieke variabelen te bepalen, wordt gebruikgemaakt van SHAP (SHapley Additive exPlanations). SHAP-waarden kwantificeren de bijdrage van elke feature aan de voorspelling van het model.

Voor elke speler wordt de gemiddelde absolute SHAP-waarde per feature berekend over meerdere wedstrijden. Deze waarden representeren de effectieve impact van fysieke prestaties op de modeloutput.

\section{Constructie van de fysieke spelersscore}

De fysieke spelersscore wordt geconstrueerd op basis van de geaggregeerde SHAP-waarden. Concreet wordt per speler:

\begin{enumerate}
    \item De score berekend per wedstrijd op basis van de gespeelde positie.
    \item Een gewogen gemiddelde berekend voor alle wedstrijden van de speler in een bepaalde competitie in een bepaald seizoen.
    \item Geschaald naar een interpreteerbare score tussen 0 en 100.
\end{enumerate}

De resulterende score weerspiegelt de relatieve fysieke bijdrage van een speler, op een reproduceerbare en data-gedreven manier.